\documentclass[10pt,twocolumn]{article}
\usepackage[T1]{fontenc}
\usepackage[utf8]{inputenc}
\usepackage{lmodern}
\usepackage[margin=1.8cm,top=2.4cm,bottom=2cm,columnsep=0.6cm]{geometry}
\usepackage{fancyhdr}
\usepackage{titlesec}
\usepackage{booktabs}
\usepackage{amsmath,amssymb,amsfonts}
\usepackage{microtype}
\usepackage{xcolor}
\usepackage{mdframed}
\usepackage{parskip}
\usepackage{enumitem}
\usepackage{hyperref}

\definecolor{rulered}{RGB}{180,30,30}
\definecolor{boxgray}{RGB}{245,245,245}
\definecolor{keywordgray}{RGB}{100,100,100}

\pagestyle{fancy}
\fancyhf{}
\fancyhead[L]{\textsc{\small Claude Mag}}
\fancyhead[C]{\small\textit{Machine Learning Research Digest}}
\fancyhead[R]{\small 2026-11-18}
\fancyfoot[C]{\small\thepage}
\renewcommand{\headrulewidth}{0.8pt}

\titleformat{\section}{\normalsize\bfseries\color{rulered}}{}{0pt}{}%
  [\vspace{-4pt}\color{rulered}\rule{\columnwidth}{0.4pt}]
\titlespacing{\section}{0pt}{8pt}{4pt}

\hypersetup{colorlinks=true,urlcolor=rulered,linkcolor=black}

\begin{document}
\setlength{\parindent}{0pt}
\setlength{\parskip}{4pt}

{\small\color{keywordgray}\textbf{Keywords:}
  Autonomous Driving \textbullet{} Masked Diffusion \textbullet{} Vision-Language Models \textbullet{} Trajectory Planning}\par\vspace{4pt}

{\LARGE\bfseries\raggedright Towards Robust Vision-Language Reasoning: A Multi-Stage Pipeline with Masked Diffusion}\par\vspace{4pt}
{\large\itshape\raggedright MD-Drive Cuts Trajectory Error by 1.68 on DriveCoT and Scores 88.9 PDMS on NAVSIM}\par\vspace{6pt}
{\small\textbf{By} Shunming Liu, Pei Huang, Yajing Sun, Huaxiao Liu}\par\vspace{2pt}
{\small\color{keywordgray}arXiv:2611.15238 \textbullet{} cs.AI, cs.CV \textbullet{} Submitted 11 Nov 2026}
\par\vspace{2pt}\noindent{\color{rulered}\rule{\columnwidth}{0.6pt}}\vspace{6pt}

Liu~et~al.\ build Masked Diffusion Drive (MD-Drive) on a three-stage VLM optimization pipeline for autonomous driving planning.
MD-Drive achieves a PDMS of 88.9 on NAVSIM, outperforming DrivingGPT by 6.5 points, and an average L2 trajectory error of 0.78 on DriveCoT---1.68 better than the strongest autoregressive baseline---representing a 29.7\% gain over the unoptimized backbone.

\begin{mdframed}[backgroundcolor=boxgray,linecolor=rulered,linewidth=1pt,
    innerleftmargin=6pt,innerrightmargin=6pt,
    innertopmargin=6pt,innerbottommargin=6pt]
\textbf{\small AT A GLANCE}\par\vspace{4pt}
\begin{description}[leftmargin=0pt,labelsep=4pt,itemsep=2pt,parsep=0pt,
    font=\small\bfseries]
  \item[Problem:] Autoregressive VLMs generate trajectories token-by-token, creating sequential dependencies and failing to capture the multi-modal distribution of valid driving futures.
  \item[Why it matters:] Planning quality directly determines safety in autonomous driving; small improvements in L2 error translate to fewer near-misses.
  \item[Method:] Three-stage VLM training (visual pre-training, driving fine-tuning, diffusion head training) followed by masked diffusion trajectory decoding.
  \item[Result:] PDMS 88.9 on NAVSIM; L2 error 0.78 on DriveCoT; 29.7\% gain over unoptimized backbone.
  \item[Evidence:] NAVSIM, DriveCoT, nuScenes benchmarks; ablations on training stages.
\end{description}
\end{mdframed}

\section{The Big Picture}

Autonomous driving planning has benefited enormously from large vision-language models: by grounding scene understanding in natural language descriptions of road conditions, weather, and traffic, VLMs can reason about complex scenarios that overwhelm purely geometric planners.

However, the generation mechanism matters.
Standard autoregressive VLMs decode trajectories token by token, left to right.
This creates two problems.
First, early tokens commit the trajectory to a direction before later tokens can correct for scene dynamics observed mid-sequence.
Second, autoregressive models generate a single trajectory mode, unable to represent the \emph{distribution} of valid futures (e.g., a highway merge can be executed earlier or later, with different spacing, all correctly).

Diffusion models address both problems by iterative denoising: a trajectory is initialized as noise and refined over $K$ steps, with each step able to attend to the entire trajectory simultaneously.
The key question for MD-Drive is whether a masked diffusion head can be integrated into a VLM backbone through a principled training pipeline---without the backbone losing its visual reasoning capabilities.

\section{Why Existing Approaches Fall Short}

\textbf{Autoregressive planners} (DrivingGPT, FSDrive) produce good scene descriptions but generate trajectories with the same token-by-token mechanism used for text, losing the multi-modality and temporal coherence that diffusion provides.

\textbf{Separate diffusion planners} operate on features from a frozen VLM encoder.
The frozen encoder was not trained to produce representations optimal for diffusion-based trajectory generation, leaving performance on the table.

\textbf{Prior discrete diffusion methods} (MVLAD-AD) introduce a diffusion planner but do not use a multi-stage optimization pipeline to align the backbone with the planner.
MD-Drive shows that this alignment is the key driver of the performance gain: ablating the pipeline while keeping the diffusion head reduces PDMS by approximately 8 points.

\section{Core Mechanism}

MD-Drive has three interacting components:
\begin{enumerate}[leftmargin=1.5em,itemsep=2pt]
  \item \textbf{VLM backbone:} A vision-language model trained through the three-stage pipeline to produce scene embeddings conditioned on the task description.
  \item \textbf{Trajectory tokenizer:} Encodes continuous trajectory waypoints $\mathbf{x}_{1:T}$ into discrete tokens via VQ-VAE.
  \item \textbf{Masked diffusion head:} Decodes trajectory tokens iteratively, conditioning on the backbone's scene embedding at each denoising step.
\end{enumerate}

\section{Technical Formulation}

\subsection*{Three-Stage Pipeline}

\textbf{Stage 1 (Visual Pre-training):}
The backbone $f_\phi$ is trained on a large-scale driving scene description dataset:
\begin{equation}
\mathcal{L}_{\text{S1}} = -\mathbb{E}\bigl[\log p_\phi(\text{caption} \mid \text{image})\bigr]
\end{equation}

\textbf{Stage 2 (Driving Fine-tuning):}
The backbone is fine-tuned on paired driving descriptions and reference trajectories:
\begin{equation}
\mathcal{L}_{\text{S2}} = -\mathbb{E}\bigl[\log p_\phi(\text{decision} \mid \text{image}, \text{scene description})\bigr]
\end{equation}

\textbf{Stage 3 (Diffusion Head):}
The backbone is frozen; the diffusion head $h_\psi$ is trained:
\begin{equation}
\mathcal{L}_{\text{S3}} = \mathbb{E}_{\mathbf{x}^{(0)}, k}\Bigl[\bigl\|\mathbf{x}^{(0)} - h_\psi(\mathbf{x}^{(k)}, k, \mathbf{z}_{\text{scene}})\bigr\|^2\Bigr]
\end{equation}
where $\mathbf{z}_{\text{scene}} = f_\phi(\text{image}, \text{description})$ and $\mathbf{x}^{(k)}$ is the trajectory at diffusion step $k$.

\subsection*{Masked Diffusion Denoising}

At diffusion step $k$, a subset of trajectory tokens $\mathbf{x}^{(k)}$ is masked:
\begin{equation}
q(\mathbf{x}^{(k)} \mid \mathbf{x}^{(0)}) = \prod_{t=1}^{T} \text{Mask}(x_t, k/K)
\end{equation}
where tokens are independently masked with probability $k/K$ (more masking at earlier steps).
The head predicts masked tokens:
\begin{equation}
p_\psi(\mathbf{x}^{(0)} \mid \mathbf{x}^{(k)}, \mathbf{z}_{\text{scene}}) = \prod_t p_\psi(x_t \mid \mathbf{x}^{(k)}_{\neq t}, \mathbf{z}_{\text{scene}})
\label{eq:denoise}
\end{equation}
At inference, $K=10$ denoising steps are used, initializing $\mathbf{x}^{(K)}$ as fully masked.

\section{Learning or Inference Procedure}

\textbf{Training:}
Stages 1 and 2 require no special hardware beyond standard VLM fine-tuning.
Stage 3 trains the diffusion head with batch size 128 for 100K steps on a 4-GPU cluster, taking approximately 24 hours.
The backbone is frozen in Stage 3 to preserve visual reasoning capabilities.

\textbf{Inference:}
Given an image and task description:
\begin{enumerate}[leftmargin=1.5em,itemsep=1pt]
  \item Compute scene embedding $\mathbf{z}_{\text{scene}} = f_\phi(\text{image}, \text{description})$.
  \item Initialize trajectory as fully masked: $\mathbf{x}^{(K)} = [\text{MASK}]^T$.
  \item Iteratively denoise using Eq.~(\ref{eq:denoise}) for $K$ steps.
  \item Decode the final trajectory $\mathbf{x}^{(0)}$ via VQ-VAE decoder.
\end{enumerate}
Total inference latency is 243.9 tokens/s, comparable to autoregressive baselines.

\section{What the Guarantee Says}

No formal planning guarantees are provided.
The empirical evidence covers three driving benchmarks with diverse conditions.
The 29.7\% gain is measured on the primary planning metric (PDMS); the L2 improvement on DriveCoT provides a geometry-based validation that is independent of the PDMS formulation.

\section{Experimental Findings}

\begin{center}
\small
\begin{tabular}{lrr}
\toprule
\textbf{Method} & \textbf{PDMS\,$\uparrow$} & \textbf{DriveCoT L2\,$\downarrow$} \\
\midrule
DrivingGPT        & 82.4 & 2.84 \\
FSDrive           & 85.1 & 2.46 \\
MVLAD-AD          & 86.3 & 1.28 \\
MD-Drive (ours)   & \textbf{88.9} & \textbf{0.78} \\
\bottomrule
\end{tabular}
\end{center}

The 29.7\% gain over the unoptimized backbone (PDMS~68.5, L2~1.10) reflects the cumulative contribution of all three training stages.
Ablating the multi-stage pipeline and training the diffusion head on a vanilla backbone reduces PDMS to approximately 80.4.

\section{Ablations and Interpretation}

\begin{itemize}[leftmargin=1em,itemsep=2pt]
  \item \textbf{Stage 1 only:} +3.5 PDMS over unoptimized backbone; establishes visual scene representations.
  \item \textbf{Stages 1--2:} +7.1 PDMS; driving-specific fine-tuning aligns representations with planning decisions.
  \item \textbf{Stages 1--2 + diffusion head (MD-Drive):} +20.4 PDMS; the full pipeline is synergistic.
  \item \textbf{Autoregressive head on Stage-1--2 backbone:} +11.3 PDMS; diffusion adds 9.1 points beyond autoregressive decoding, confirming the value of multi-modal trajectory generation.
  \item \textbf{Denoising steps $K$:} $K=10$ is optimal; $K=5$ degrades L2 by 0.12; $K=20$ provides no additional benefit.
\end{itemize}

\section*{Reference}
S.\ Liu, P.\ Huang, Y.\ Sun, H.\ Liu.
\textbf{Towards Robust Vision-Language Reasoning: A Multi-Stage Optimization Pipeline with Masked Diffusion Model.}
arXiv:2611.15238. cs.AI, cs.CV. November 2026.
\url{https://arxiv.org/abs/2611.15238}

\end{document}
